% !Bibliography = biber
\documentclass[12pt,a4paper,oneside]{report}

\input{preambule}

\begin{document}

\pagenumbering{roman}

\begin{titlepage}
    \centering
    {\Large Université Lyon 2, Lyon 3, ENSIB, ENS Lyon\par}
    \vspace{2cm}
    {\Huge\bfseries Identifier les marqueurs du naturalisme par la modélisation thématique et la classification automatique : signature d’auteur, contenu thématique et biais de corpus\par}
    \vspace{2cm}
    {\Large Mémoire présenté par Morgan RICHARD\par}
    \vfill
    {\large Année universitaire 2025--2026\par}
\end{titlepage}

\chapter*{Remerciements}
\addcontentsline{toc}{chapter}{Remerciements}

Ce mémoire a été réalisé dans le cadre de mon stage de fin d’études au sein du laboratoire LATTICE, unité mixte de recherche du CNRS, de l’École Normale Supérieure-PSL et de l’Université Sorbonne Nouvelle. Les recherches du laboratoire portent notamment sur la linguistique, le traitement automatique des langues, l’évolution linguistique et les humanités numériques. Ce cadre scientifique m’a permis de mobiliser des méthodes de modélisation thématique et de classification automatique afin d’étudier les régularités lexicales et thématiques associées à l’écriture naturaliste. Le stage a été encadré par Dominique Legallois que je remercie pour son encadrement et pour toutes les nouvelles notions qu'il m'a enseignées au long de ce stage. Je tiens à remercier toutes les autres personnes du laboratoire pour leur acceuil chaleureux et leur soutien tout au long de ce projet, notament à Virgie Pauchont, gestionnaire et secrétaire, à Yoann Dupont, MCF à Sorbonne Nouvelle, à Mathieu Dehouck Chargé de Recherche au CNRS, et à Olga Seminck, ingénieure de recherche au CNRS.

Je remercie également Madame Pegah pour son encadrement ainsi que tous le personnel encadrant du master HN de Lyon pour leurs conseils tout au long de l'année.

Je remercie également mes proches pour leur soutien et leurs encouragements, je pense notamment à mes parents ainsi qu'à ma copine. 

\chapter*{Résumé}
\addcontentsline{toc}{chapter}{Résumé}

Ce mémoire porte sur...

\tableofcontents
\listoffigures
\listoftables

\clearpage
\pagenumbering{arabic}

\chapter*{Introduction}
\addcontentsline{toc}{chapter}{Introduction}

L'objectif de ce mémoire est de présenter les travaux réalisés dans le cadre de mon stage de fin d'études, effectué au sein du laboratoire de recherche LATTICE\footnote{Le Lattice est un laboratoire du CNRS (UMR8094), rattaché à l’Ecole Normale Supérieure (IDEX PSL : Paris Sciences et Lettres) et à l’Université Sorbonne Nouvelle-Paris 3. Il regroupe des chercheurs CNRS, des enseignants-chercheurs de l’Université Paris Sorbonne Nouvelle et d’autres universités, ainsi que de nombreux autres personnels permanents ou non-permanents.
Les recherches du Lattice sont consacrées à la linguistique et au traitement automatique des langues.Dans le domaine linguistique proprement dit, nos recherches portent, d’une part, sur l’interface entre grammaire et lexique (dans le courant des grammaires cognitives), et d’autre part, sur le discours (en particulier la structuration du discours et la coréférence).
Le Lattice accorde également une place majeure à l’étude du changement linguistique, qu’il s’agisse de l’évolution du français – en particulier dans les domaines syntaxique et sémantique – ou des tendances plus générales à l’œuvre dans l’évolution des langues.
Dans le domaine du TAL, le Lattice explore les techniques neuronales depuis plusieurs années, et a récemment mis l’accent sur l’analyse syntaxique. Les travaux en linguistique et en TAL sont étroitement corrélés.
Enfin notre laboratoire est fortement impliqué dans le domaine des Humanités Numériques et contribue à la construction de ressources enrichies. https://www.lattice.cnrs.fr/ }. Au cours de ce stage, j’ai travaillé sur de nombreux aspects de \textit{TAL} (Traitement Automatique des Langues). Mon objectif principal était de déterminer comment était redigée l'ecriture naturaliste. Ce stage a été encadré par le professeur Dominque Legallois ainsi que par d'autre chercheurs du laboratoire à qui j'ai pu solicité leur aide notament lors d'utilisation de technique de \textit{machine Learning}.

La modélisation thématique, la stylométrie et la classification automatique constituent trois approches complémentaires de l’analyse quantitative des textes. Elles se distinguent principalement par leurs objectifs : la modélisation thématique cherche à faire émerger les structures sémantiques d’un corpus, la stylométrie mesure les caractéristiques formelles de l’écriture, tandis que la classification automatique attribue aux textes des catégories préalablement définies.

La modélisation thématique (\emph{topic modeling}) vise généralement à dégager, sans annotation préalable, les principaux thèmes présents dans un ensemble de documents. L’allocation de Dirichlet latente (\emph{Latent Dirichlet Allocation}, LDA) s’est imposée comme l’un des modèles probabilistes de référence. Elle représente chaque document comme un mélange de thèmes, eux-mêmes définis par des distributions de probabilités sur les mots du corpus.\footnote{David M. Blei, Andrew Y. Ng et Michael I. Jordan, « Latent Dirichlet Allocation », \emph{Journal of Machine Learning Research}, vol.~3, 2003, p.~993--1022, \href{https://www.jmlr.org/papers/v3/blei03a.html}{texte intégral en ligne}.} Les prolongements de cette méthode permettent notamment d’intégrer l’évolution temporelle des thèmes, leurs relations ou les métadonnées associées aux documents.\footnote{David M. Blei, « Probabilistic Topic Models », \emph{Communications of the ACM}, vol.~55, no~4, 2012, p.~77--84, \href{https://www.cs.columbia.edu/~blei/papers/Blei2012.pdf}{texte intégral en ligne}.} Plus récemment, les modèles neuronaux et les plongements contextuels ont renouvelé la modélisation thématique, en particulier pour l’analyse des textes courts, multilingues ou sémantiquement complexes.\footnote{Xiaobao Wu, Thong Nguyen et Anh Tuan Luu, « A Survey on Neural Topic Models: Methods, Applications, and Challenges », \emph{Artificial Intelligence Review}, vol.~57, article~18, 2024, \href{https://doi.org/10.1007/s10462-023-10661-7}{DOI et texte intégral}.} BERTopic, par exemple, combine des représentations produites par un modèle de langue, une méthode de regroupement automatique et une variante du TF--IDF pour caractériser les thèmes obtenus.\footnote{Maarten Grootendorst, « BERTopic: Neural Topic Modeling with a Class-Based TF-IDF Procedure », prépublication arXiv, 2022, \href{https://arxiv.org/abs/2203.05794}{texte intégral en ligne}.} L’interprétation des résultats demeure néanmoins délicate : les thèmes dépendent de la constitution et du prétraitement du corpus, du nombre de thèmes demandé et des paramètres du modèle. Par ailleurs, une bonne performance statistique ne garantit pas nécessairement que les thèmes produits soient cohérents pour un lecteur humain.\footnote{Jonathan Chang, Jordan Boyd-Graber, Sean Gerrish, Chong Wang et David M. Blei, « Reading Tea Leaves: How Humans Interpret Topic Models », dans \emph{Advances in Neural Information Processing Systems 22}, 2009, p.~288--296, \href{https://papers.neurips.cc/paper/2009/hash/f92586a25bb3145facd64ab20fd554ff-Abstract.html}{texte intégral en ligne}.}

La stylométrie cherche, pour sa part, à mesurer les régularités formelles d’une écriture à partir de caractéristiques quantifiables, telles que les fréquences lexicales, l’emploi des mots-outils, la richesse du vocabulaire, la longueur des phrases, la ponctuation, les catégories grammaticales ou les séquences de caractères. Elle a d’abord été principalement utilisée pour résoudre des problèmes d’attribution d’auteur, comme l’illustre l’étude fondatrice de Frederick Mosteller et David Wallace sur les articles anonymes des \emph{Federalist Papers}.\footnote{Frederick Mosteller et David L. Wallace, \emph{Inference and Disputed Authorship: The Federalist}, Reading, Addison-Wesley, 1964 ; réédition avec une introduction de John Nerbonne, Stanford, CSLI Publications, 2007, \href{https://web.stanford.edu/group/cslipublications/cslipublications/site/1575865521.shtml}{notice éditoriale}.} Parmi les méthodes classiques, le \emph{Delta} de John Burrows mesure la distance stylistique entre plusieurs textes à partir des fréquences standardisées de leurs mots les plus courants.\footnote{John Burrows, « `Delta': A Measure of Stylistic Difference and a Guide to Likely Authorship », \emph{Literary and Linguistic Computing}, vol.~17, no~3, 2002, p.~267--287, \href{https://doi.org/10.1093/llc/17.3.267}{DOI}.} Les recherches contemporaines mobilisent également les n-grammes de caractères, les méthodes d’apprentissage automatique et les représentations produites par les modèles de langue.\footnote{Efstathios Stamatatos, « A Survey of Modern Authorship Attribution Methods », \emph{Journal of the American Society for Information Science and Technology}, vol.~60, no~3, 2009, p.~538--556, \href{https://www.clips.uantwerpen.be/~walter/educational/material/Stamatatos_survey2009.pdf}{texte intégral en ligne}.} La stylométrie comprend désormais plusieurs tâches, notamment l’attribution, la vérification et le profilage d’auteur, ainsi que l’étude de l’évolution du style au cours du temps.\footnote{Enzo Terreau, \emph{Apprentissage de représentations d’auteurs et d’autrices à partir de modèles de langue pour l’analyse des dynamiques d’écriture}, thèse de doctorat en informatique, Université Lyon~2, 2024, \href{https://theses.fr/2024LYO20001}{notice et résumé en ligne}.} Ses résultats restent toutefois sensibles à la longueur des documents et aux effets du genre, du thème, de la période historique, du support ou de la traduction. Ces facteurs doivent donc être contrôlés afin de ne pas les confondre avec le style individuel.

Enfin, la classification automatique consiste à attribuer un texte à une ou plusieurs catégories définies à l’avance. Contrairement à la modélisation thématique, elle relève généralement de l’apprentissage supervisé : le modèle apprend à partir d’un corpus dont les documents ont préalablement été annotés.\footnote{Fabrizio Sebastiani, « Machine Learning in Automated Text Categorization », \emph{ACM Computing Surveys}, vol.~34, no~1, 2002, p.~1--47, \href{https://arxiv.org/abs/cs/0110053}{texte intégral en ligne}.} Les approches classiques représentent les documents par des fréquences lexicales ou des pondérations TF--IDF, puis utilisent des méthodes comme le classifieur bayésien naïf, la régression logistique ou les machines à vecteurs de support.\footnote{Christopher D. Manning, Prabhakar Raghavan et Hinrich Schütze, \emph{Introduction to Information Retrieval}, Cambridge, Cambridge University Press, 2008, chap.~13 à 15, \href{https://nlp.stanford.edu/IR-book/information-retrieval-book.html}{édition intégrale en ligne}.} Ces méthodes ont été complétées par les réseaux de neurones, puis par les modèles de langue préentraînés fondés sur l’architecture Transformer. BERT a notamment popularisé une procédure dans laquelle un modèle préentraîné sur de grandes quantités de textes est ensuite spécialisé sur une tâche de classification à partir d’un corpus annoté plus restreint.\footnote{Jacob Devlin, Ming-Wei Chang, Kenton Lee et Kristina Toutanova, « BERT: Pre-training of Deep Bidirectional Transformers for Language Understanding », dans \emph{Proceedings of NAACL-HLT 2019}, p.~4171--4186, \href{https://aclanthology.org/N19-1423/}{texte intégral en ligne}.} Les grands modèles de langue permettent également d’effectuer certaines classifications sans apprentissage spécifique ou à partir d’un nombre limité d’exemples. Ils ne rendent cependant pas systématiquement obsolètes les modèles spécialisés : pour de nombreuses tâches, un modèle plus petit correctement entraîné peut rester plus performant, plus stable et moins coûteux.\footnote{Aleksandra Edwards et Jose Camacho-Collados, « Language Models for Text Classification: Is In-Context Learning Enough? », dans \emph{Proceedings of LREC-COLING 2024}, p.~10058--10072, \href{https://aclanthology.org/2024.lrec-main.879/}{texte intégral en ligne}.}

Dans ces trois domaines, le choix de la méthode ne peut être séparé des conditions de constitution et d’annotation du corpus. Une évaluation rigoureuse suppose notamment de contrôler les effets d’auteur, de source, de genre et de période, de séparer strictement les données d’apprentissage et de test, et de confronter les résultats quantitatifs à une lecture qualitative des textes.\footnote{Justin Grimmer et Brandon M. Stewart, « Text as Data: The Promise and Pitfalls of Automatic Content Analysis Methods for Political Texts », \emph{Political Analysis}, vol.~21, no~3, 2013, p.~267--297, \href{https://doi.org/10.1093/pan/mps028}{DOI}.}
 
Mon obje
Le naturalisme constitue l’un des principaux mouvements littéraires et artistiques de la seconde moitié du XIXe siècle. Prolongeant certaines ambitions du réalisme, il se donne pour objectif de représenter l’être humain et la société avec une précision inspirée des sciences expérimentales. L’œuvre littéraire ne repose donc plus seulement sur l’imagination de l’écrivain : elle s’appuie sur l’observation du réel, sur une documentation minutieuse ainsi que sur l’étude des mécanismes physiologiques, sociaux et psychologiques qui déterminent les comportements individuels.

Émile Zola\footnote{Pour ce qui est des informations sur Émile Zola, je me suis grandement référé à sa page Wikipédia, qui est considérée comme une page de qualité : \href{{https://fr.wikipedia.org/wiki/\%C3\%89mile_Zola}}{lien de sa page.}}, principal théoricien du mouvement, définit ainsi le naturalisme comme « l’esprit scientifique porté dans toutes nos connaissances ». Il entend substituer à « l’homme métaphysique » un « homme physiologique », inséparable de son hérédité et du milieu dans lequel il évolue\footnote{\url{https://fr.wikipedia.org/wiki/Naturalisme_(litt\%C3\%A9rature)}}. Inspiré notamment par \textit{l’Introduction à l’étude de la médecine expérimentale} de Claude Bernard, Zola conçoit le romancier comme un observateur et un expérimentateur. Celui-ci commence par recueillir et consigner des faits, puis place ses personnages dans des circonstances déterminées afin d’étudier le fonctionnement de leurs passions, de leurs instincts et de leurs conduites. Le roman devient alors une forme d’expérience permettant de mettre au jour les lois qui régissent l’individu et la société\footnote{Émile Zola, \textit{Le Roman expérimental}, Paris, Charpentier, 1881.}.

Cette conception repose principalement sur le déterminisme. Dans la perspective naturaliste, la personnalité et la destinée d’un individu sont largement conditionnées par son héritage physiologique, son environnement familial, sa situation économique et son appartenance sociale. Le naturalisme ne doit cependant pas être réduit à un ensemble de procédés stylistiques\footnote{\url{https://fr.wikipedia.org/wiki/Naturalisme_(litt\%C3\%A9rature)}}. Paul Alexis le présente avant tout comme une méthode permettant de penser\footnote{Jules Huret, \textit{Enquête sur l'évolution littéraire}, Paris, Charpentier, 1891.}, de voir, d’étudier et d’expérimenter. De même, Camille Lemonnier l’envisage comme un « réalisme agrandi » par l’analyse des milieux et par l’apport des sciences naturelles et sociales. Le mouvement s’inscrit ainsi dans une vision matérialiste et empirique du monde, selon laquelle les phénomènes humains doivent être expliqués à partir de causes observables\footnote{\url{https://fr.wikipedia.org/wiki/Naturalisme_(litt\%C3\%A9rature)}}.

Sur le plan esthétique et narratif, cette ambition scientifique entraîne une remise en cause des conventions romanesques traditionnelles. Les naturalistes rejettent les intrigues artificielles, les péripéties invraisemblables et les personnages idéalisés. À la construction d’une fable spectaculaire, ils préfèrent la restitution d’une « tranche de vie », c’est-à-dire l’étude détaillée d’un fragment d’existence ordinaire. Le roman prend alors la forme d’une monographie consacrée à un individu, à une famille, à un métier ou à un milieu social. Le narrateur tend à s’effacer derrière les faits et adopte une posture d’impersonnalité comparable à celle d’un greffier chargé d’exposer les pièces d’un dossier sans formuler explicitement de jugement\footnote{Émile Zola, \textit{Le Roman expérimental}}.

La description occupe, dans cette perspective, une fonction essentielle. Elle ne constitue pas un simple ornement esthétique, mais un instrument d’analyse et d’explication. Les logements, les vêtements, les paysages, les lieux de travail et les conditions matérielles d’existence contribuent à façonner les personnages. Décrire un environnement revient donc à rendre visibles les forces qui agissent sur l’individu. Cette attention portée aux déterminations sociales conduit les écrivains naturalistes à s’intéresser aux catégories jusque-là marginalisées par la littérature : les ouvriers, les paysans, les prostituées, les criminels, les malades ou les exclus. Le héros traditionnel s’efface au profit de personnages ordinaires, souvent confrontés à la pauvreté, à la violence, à la maladie ou au déclassement.

La représentation de ces réalités a toutefois suscité d’importantes controverses morales et esthétiques. La crudité de certaines scènes, l’étude des pathologies et l’attention accordée aux formes de déchéance ont conduit les adversaires du mouvement à l’accuser de complaisance envers la laideur et le vice. Zola défend au contraire la dimension morale de sa démarche. À ses yeux, le romancier agit comme un anatomiste ou un chirurgien : en analysant les causes matérielles et sociales des comportements, il contribue à mieux comprendre les dysfonctionnements de la société et fournit les connaissances nécessaires à leur éventuelle correction. L’écrivain naturaliste serait ainsi un « moraliste expérimentateur », dont la fonction n’est pas de condamner, mais d’expliquer\footnote{\textit{Ibid.}}.

Cette conception est notamment contestée par Ferdinand Brunetière\footnote{Ferdinand Brunetière, \textit{Le roman naturaliste}, Paris, Calmann-Lévy, 1883.}, qui reproche au naturalisme français de privilégier la recherche de la grossièreté, le pessimisme et l’« écriture artiste » au détriment de la profondeur psychologique. Il oppose à ce qu’il considère comme un « faux naturalisme » une littérature incarnée par Dickens, George Eliot ou Tolstoï, caractérisée par la probité de l’observation, la simplicité de l’exécution et la sympathie envers les humbles. Edmond Cattier insiste également sur le rôle de l’« imagination sympathique »\footnote{Edmond Cattier, \textit{Le Naturalisme littéraire}, Bruxelles, Weissenbruch, 1897.}, faculté par laquelle l’écrivain dépasse la simple observation extérieure pour pénétrer la conscience d’autrui, comprendre les motifs de ses actions et restituer la complexité de son expérience. Ces critiques mettent en évidence une tension centrale du naturalisme : celle qui oppose l’ambition d’objectivité scientifique à la nécessité d’une compréhension sensible et psychologique de l’être humain.

Le naturalisme ne forme d’ailleurs pas une école parfaitement homogène. S’il revendique l’héritage de Diderot, de Rousseau, de Balzac, de Stendhal, de Flaubert et des frères Goncourt, ses représentants divergent sur la définition même du mouvement. Edmond de Goncourt préfère parler de « naturisme », Joris-Karl Huysmans d’« intimisme », Louis-Marie Desprez d’« impressionnisme », tandis que Guy de Maupassant propose le terme d’« illusionnistes ». Cette diversité terminologique révèle l’existence de pratiques et de conceptions distinctes derrière l’apparente unité du groupe naturaliste\footnote{\url{https://fr.wikipedia.org/wiki/Naturalisme_(litt\%C3\%A9rature)}}.

Le succès des \textit{Soirées de Médan}, publiées en 1880, marque l’un des moments les plus emblématiques de cette aventure collective. Toutefois, l’unité du mouvement se fragilise rapidement. Le \textit{Manifeste des Cinq}, publié en 1887 contre \textit{La Terre} de Zola, rend visibles les désaccords internes, tandis que plusieurs auteurs s’éloignent progressivement de l’esthétique naturaliste. Huysmans, notamment, ouvre avec \textit{À rebours} la voie à une littérature décadente et symboliste\footnote{\textit{Ibid.}}. À la fin du siècle, le développement du roman psychologique et du symbolisme contribue ainsi au déclin du naturalisme en tant qu’école constituée. La formule attribuée à Stéphane Mallarmé, selon laquelle « le mot mourra quand Zola aura achevé son œuvre », souligne à la fois le rôle déterminant de Zola dans la définition du mouvement et la difficulté du naturalisme à lui survivre sous une forme unifiée\footnote{Jules Huret, \textit{Enquête sur l'évolution littéraire}, Paris, Charpentier, 1891.}.

Dans quelle mesure les méthodes de modélisation thématique et de classification automatique permettent-elles d’identifier des marqueurs du naturalisme sans confondre signature d’auteur, contenu thématique et artefacts de constitution du corpus ?

Afin de répondre à cette problématique, ce mémoire s’organise en trois chapitres complémentaires. Le premier présente une exploration thématique classique des trente et un romans d’Émile Zola à l’aide de la factorisation en matrices non négatives et de la méthode K-means. Cette première analyse vise à dégager des régularités lexicales et thématiques internes à l’œuvre zolienne, considérées comme des marqueurs candidats plutôt que comme des preuves immédiates du naturalisme. Le deuxième chapitre étudie ces régularités au moyen de BERTopic et analyse leur évolution au cours de la carrière de Zola. Le troisième chapitre confronte les hypothèses issues de ces analyses à un corpus comparatif grâce à une classification automatique opposant romans naturalistes et romans relevant d’autres courants. Une attention particulière est accordée aux risques de confusion entre signature d’auteur, contenu thématique et artefacts liés à la constitution du corpus. Enfin, les choix de nettoyage, l’organisation de la base de données et les principaux éléments nécessaires à la reproductibilité des expériences sont présentés en annexe.






\input{chapitres/chapitre_1}
\chapter{Dynamique Topic modeling sur le corpus de d'Émile Zola}

\section{présentation des données}

\subsection{Présentation générale}

Ici, je vais presenter une autres façon de faire du topic modeling sur le corpus d'Émile Zola. Cette fois-ci, je vais utiliser une approche de \textit{Dynamic Topic Modeling} (DTM) pour analyser l'évolution des thèmes au fil du temps dans les romans de Zola. L'objectif est de comprendre comment les motifs naturalistes se transforment et circulent à travers son œuvre, du début de sa carrière jusqu'à ses derniers romans.

Le corpus que j'ai utlisé pour faire le DTM est le même\footnote{Cf. chapitre 1 dans la sous partie présentation des données} que celui utilisé pour le topic modéling classique\footnote{Ici quand je dis \textit{Topic Modeling} classique, c'est du \textit{Topic Modeling} statique ou figé sur l'ensemble global de son oeuvre}. 

Ici ce qui sera intéressent, c'est de voir comment les topics évoluent au fil du temps, voir quels topics apparaissent, disparaissent, se transforment ou plus intéressent voir quels topic restent constant dans le temps au cours de la carrière de Zola. Voir quels topics restent constant dans le temps permettra de voir si certains thèmes que l'ont sait naturalistes sont bien présent dans l'ensemble de son oeuvre ou si certains apparaissent seulement à un moment donné de sa carrière.

Pour ce qui est de la méthodologie, j'ai utilisé le même prétraitement que pour le topic modeling classique. J'ai donc tokenisé les textes, supprimé les stop words, et appliqué une lemmatisation pour réduire les mots à leur forme de base. 

Les données sont organisées de manière à ce que chaque roman soit associé a une année de publication, ce qui permet de suivre l'évolution des topics au fil du temps. 

\subsection{Statistique exploratoires}

Les romans sont segementés en paquets de 40 phrases, c'est la même segmentation que pour le \textit{Topic Modeling} dans le chapitre 1, cela permet d'avoir une meilleure granularité dans l'analyse des topics. 

On se retrouve donc avec un jeu de donnée 5327 documents\footnote{C'est à dire 5327 paquets de phrase}.

\begin{table}[htbp]
\centering
\caption{Statistiques descriptives de la variable \textit{phrases\_paquet}}
\label{tab:stat_phrases_paquet}
\begin{tabular}{lr}
\toprule
\textbf{Métrique} & \textbf{Valeur} \\
\midrule
count & 5327.000 \\
mean  & 785.298 \\
std   & 212.129 \\
min   & 42.000 \\
25\%  & 636.000 \\
50\%  & 753.000 \\
75\%  & 902.500 \\
max   & 2944.000 \\
\bottomrule
\end{tabular}
\end{table}

En moyenne, un paquet contient environ 785 mots, tandis que la médiane s’élève à 753 mots, ce qui indique que la distribution est légèrement influencée par certains paquets particulièrement longs. La moitié des documents comprend entre 636 et 902,5 mots, ce qui montre que la majorité des paquets possède une taille relativement homogène malgré un écart type de 212 mots. Les valeurs extrêmes, comprises entre 42 et 2,944 mots, révèlent néanmoins une forte variabilité ponctuelle, qui peut s’expliquer par les différences de longueur des phrases entre les romans, mais aussi par la présence éventuelle de paquets incomplets, notamment à la fin de chaque œuvre. Dans l’ensemble, ces résultats confirment que la segmentation en 40 phrases produit des unités textuelles suffisamment riches pour l’analyse thématique, tout en conservant une certaine diversité de longueur.

L’examen des valeurs extrêmes\footnote{Le code est disponible en annaxe du document} montre que la forte amplitude observée entre le minimum et le maximum concerne en réalité une très faible proportion du corpus. Parmi les 5,327 paquets, seuls 12 contiennent moins de 300 mots, soit environ 0,23\% des documents, tandis que 24 dépassent 1,500 mots, soit environ 0,45\%. Parmi ces derniers, un seul paquet contient plus de 2,000 mots, ce qui confirme que la valeur maximale de 2,944 mots constitue un cas particulièrement isolé. Au total, 5,291 paquets, soit environ 99,32\% du jeu de données, comportent entre 300 et 1,500 mots. La variabilité mise en évidence par l’écart type est donc principalement liée à quelques observations atypiques et ne remet pas en cause l’homogénéité globale de la segmentation. Les paquets les plus courts peuvent notamment correspondre aux dernières sections des romans, lorsqu’il reste moins de 40 phrases à regrouper, tandis que les paquets les plus longs peuvent s’expliquer par la présence de phrases particulièrement développées ou par des spécificités de ponctuation et de segmentation propres à certains textes.
\section{Utilisation de BerTopic}

\subsection{Lemmatisation des données}

Comme pour tout modèle de traitement du langage naturel, il est important de préparer les données avant de les utiliser avec BerTopic. La lemmatisation est une étape clé dans ce processus, j'ai suivi les memes etapes que pour le NMF ainsi que le KMeans. 

J'ai utilisé la bibliothèque spaCy pour effectuer la lemmatisation des textes, j'ai réutilisé le même modele de langue que pour le \textit{Topic modeling} classique, le \textit{"fr\_core\_news\_lg"}. Cette étape permet de réduire les mots à leur forme de base, ce qui facilite l'analyse et l'extraction des thèmes. 

\subsection{Préparation des embeddings}

Le modèle \texttt{dangvantuan/sentence-camembert-base} est utilisé afin de transformer chaque segment textuel en un vecteur numérique dense, appelé \textit{embedding}, qui représente son contenu sémantique. Fondé sur CamemBERT et adapté à la langue française, ce modèle permet de rapprocher les documents qui traitent de thèmes similaires, même lorsqu’ils n’emploient pas exactement le même vocabulaire. Ces représentations vectorielles sont ensuite fournies à BERTopic, qui s’appuie sur leur proximité pour réduire leur dimension, regrouper les documents sémantiquement proches et identifier les principaux topics du corpus. L’utilisation de ces embeddings est particulièrement pertinente dans le cadre d’un \textit{Dynamic Topic Modeling} (DTM), puisqu’elle permet d’étudier l’évolution des thèmes au cours du récit tout en tenant compte du contexte lexical et sémantique de chaque paquet de phrases.

\subsection{Préparation des paramètres du modèle}

Afin d’identifier des thèmes cohérents à partir des représentations sémantiques, BERTopic combine une réduction de dimension avec un algorithme de classification non supervisée. Le modèle UMAP réduit les embeddings à trois dimensions afin de faciliter la détection des groupes tout en conservant au mieux leur structure sémantique. Le paramètre \texttt{n\_neighbors=20} permet de prendre en compte un voisinage suffisamment large pour préserver à la fois les relations locales et l’organisation générale du corpus, tandis que \texttt{min\_dist=0.0} favorise la formation de groupes compacts. L’utilisation de la distance cosinus est adaptée aux embeddings textuels, car elle mesure la proximité entre leurs orientations plutôt que leur amplitude. Enfin, la fixation de \texttt{random\_state=42} garantit la reproductibilité de la réduction dimensionnelle.

La classification des segments est ensuite réalisée avec HDBSCAN. Le paramètre \texttt{min\_cluster\_size=27} impose qu’un topic soit représenté par au moins 27 segments, ce qui limite l’apparition de thèmes trop spécifiques ou insuffisamment documentés. Le paramètre \texttt{min\_samples=4} conserve une certaine souplesse dans la formation des clusters, tout en permettant d’identifier comme atypiques les observations éloignées des groupes principaux. La méthode \texttt{eom} privilégie les clusters les plus stables dans la structure hiérarchique produite par HDBSCAN. L’activation de \texttt{prediction\_data} conserve par ailleurs les informations nécessaires à l’affectation ultérieure de nouveaux documents ou au traitement des observations considérées comme aberrantes.

La représentation lexicale de chaque topic repose sur un \texttt{CountVectorizer}. Le seuil \texttt{min\_df=2} élimine les termes n’apparaissant que dans un seul segment, souvent peu représentatifs ou associés à des erreurs de transcription, tandis que \texttt{max\_df=0.6} écarte les mots présents dans plus de 60\,\% des segments, dont le pouvoir discriminant est généralement limité. Le modèle \texttt{ClassTfidfTransformer} pondère ensuite les termes selon leur importance dans chaque topic. L’option \texttt{reduce\_frequent\_words=True} réduit l’influence des mots fréquents, tandis que la pondération BM25 améliore la distinction entre les termes réellement caractéristiques d’un thème et ceux qui sont simplement courants dans le corpus.
Le modèle BERTopic est enfin ajusté sur les segments lemmatisés en utilisant les embeddings précédemment calculés. La désactivation du calcul des probabilités réduit le coût computationnel, celles-ci n’étant pas indispensables à l’identification et à l’interprétation des topics. Après l’entraînement, les segments initialement classés comme observations aberrantes sont réaffectés au topic sémantiquement le plus proche à partir de leurs embeddings. Cette étape de réduction des \textit{outliers}, suivie d’une mise à jour des représentations thématiques, permet d’obtenir une partition plus complète et des topics lexicalement plus cohérents, ce qui facilite ensuite l’étude de leur évolution dans le cadre du \textit{Dynamic Topic Modeling}.
\section{resultats du Dynamic Topic Modeling sur le corpus de Zola}

\subsection{Validation mathématique}

Avant de présenter les résultats du Dynamic Topic Modeling sur le corpus de Zola, il est important de valider mathématiquement la pertinence des topics identifiés. Pour ce faire, j'ai utilisé deux méthodes complémentaires : Le score de cohérence et la silhouette.
\chapter{Machine learning sur le corpus de d'Émile Zola et auteurs contemporains}
\section{Présentation du protocole expérimental}
\label{sec:protocole_classification}

Les deux premiers chapitres ont fait apparaître plusieurs régularités lexicales et thématiques dans l'œuvre romanesque de Zola. Ces motifs constituent des marqueurs candidats, mais ils ne permettent pas encore de déterminer ce qui relève du naturalisme, de la signature personnelle de l'écrivain ou du contenu particulier de ses romans. Le présent chapitre déplace donc l'analyse vers un corpus comparatif et vers une tâche de classification supervisée. Il s'agit de déterminer si une signature textuelle apprise à partir de Zola se retrouve dans des œuvres d'auteurs absents de l'apprentissage, et si le modèle attribue des scores plus élevés aux romans étiquetés naturalistes qu'aux romans rattachés à d'autres courants.

Deux notebooks mettent cette hypothèse à l'épreuve. Le premier conserve les noms propres et correspond à la condition ANP, « avec noms propres »\footnote{Notebook ANP : \path{zola_vs_naturaliste_ANP.ipynb}.}. Le second reprend la même expérience après leur suppression automatique ; il correspond à la condition SNP, « sans noms propres »\footnote{Notebook SNP : \path{zola_vs_naturaliste_SNP.ipynb}.}. Les œuvres retenues, leur répartition, leur segmentation, leur représentation et les modèles comparés sont identiques. Seul le traitement des noms propres varie. La comparaison des deux conditions constitue ainsi une expérience d'ablation contrôlée, destinée à mesurer la part de ces indices dans les décisions du classifieur.

\subsection{Question de recherche et unité de classification}

Le protocole ne cherche pas à répartir définitivement les écrivains entre deux groupes, l'un naturaliste et l'autre non naturaliste. Une telle opposition serait peu adaptée aux trajectoires étudiées : plusieurs auteurs se sont rapprochés du naturalisme à un moment de leur carrière avant de s'en éloigner, tandis que certaines œuvres occupent une position intermédiaire. Les étiquettes « naturaliste » et « autre courant » sont donc attribuées au niveau de l'œuvre. Chacun des auteurs utilisés pour la validation et le test est représenté dans les deux classes.

La question effectivement posée est plus restreinte : un classifieur dont les seuls exemples naturalistes d'apprentissage sont les œuvres de Zola peut-il généraliser à des romans écrits par d'autres auteurs ? Une réponse positive indiquerait qu'une partie du signal appris dépasse la simple reconnaissance d'un texte zolien. Elle ne suffirait toutefois pas à établir que le modèle a isolé une définition générale du naturalisme, puisque ses exemples positifs proviennent tous d'un même écrivain. Le terme de « transfert » désigne ici cette généralisation inter-auteurs d'un signal textuel, et non le transfert d'un modèle préentraîné au sens habituel de l'apprentissage automatique.

\subsection{Constitution du corpus et séparation des auteurs}

Le fichier de métadonnées réunit initialement 191 œuvres. Trois textes ont été écartés : \textit{L'Âme étrangère} et \textit{Angelus} de Guy de Maupassant, deux fragments posthumes et inachevés, ainsi que \textit{Un dilemme} de Joris-Karl Huysmans, retenu comme une nouvelle plutôt que comme un roman. Après ces exclusions, le corpus disponible comprend 188 œuvres, dont 68 étiquetées naturalistes et 120 rattachées à d'autres courants. Les règles de partition en retiennent 107 pour l'expérience ; les 81 autres restent hors des ensembles d'apprentissage, de validation et de test.

La séparation est effectuée au niveau des auteurs avant tout découpage des textes. L'ensemble d'apprentissage contient les 31 œuvres de Zola, qui forment l'intégralité de la classe positive, et 34 œuvres appartenant à l'autre classe. Ces dernières proviennent de vingt auteurs dont toutes les œuvres retenues sont étiquetées « autre courant » et sont publiées entre 1865 et 1903, soit la période couverte par le corpus zolien. L'apprentissage réunit ainsi 65 œuvres de 21 auteurs.

L'ensemble de validation rassemble Guy de Maupassant, Camille Lemonnier et Georges Eekhoud. Il comprend 17 œuvres, dont six naturalistes et onze relevant d'autres courants. Le test réunit les frères Goncourt, Joris-Karl Huysmans et Octave Mirbeau, pour un total de 25 œuvres presque équilibré entre les deux classes. Cette construction permet d'évaluer les différences entre œuvres chez un même auteur, plutôt que de confondre systématiquement le courant littéraire et l'identité de l'écrivain.

\begin{table}[H]
    \centering
    \caption{Composition des ensembles utilisés pour la classification}
    \label{tab:composition_ensembles_classification}
    \begin{tabularx}{\textwidth}{>{\raggedright\arraybackslash}X
        >{\centering\arraybackslash}p{2cm}
        >{\centering\arraybackslash}p{2.7cm}
        >{\centering\arraybackslash}p{2.7cm}
        >{\centering\arraybackslash}p{1.8cm}}
        \toprule
        \textbf{Ensemble} & \textbf{Auteurs} & \textbf{Œuvres naturalistes} & \textbf{Autres œuvres} & \textbf{Total} \\
        \midrule
        Apprentissage & 21 & 31 & 34 & 65 \\
        Validation & 3 & 6 & 11 & 17 \\
        Test & 3 & 13 & 12 & 25 \\
        \bottomrule
    \end{tabularx}
\end{table}

Aucun auteur n'est partagé entre les trois ensembles. Le code vérifie également qu'un même texte nettoyé ne réapparaît pas sous deux titres différents, à partir de son empreinte SHA-256. Cette précaution repère les copies exactes, sans garantir l'absence de deux éditions seulement partielles ou très proches. La validation croisée aléatoire par œuvre n'a pas été retenue comme mesure principale : elle aurait réparti plusieurs romans de Zola entre différents plis et aurait surtout mesuré la reconnaissance d'un auteur déjà connu. La sélection repose au contraire sur trois auteurs entièrement absents de l'apprentissage, avant une évaluation sur trois autres auteurs également tenus à l'écart.

\subsection{Préparation et segmentation des textes}

Les textes sont normalisés selon la forme Unicode NFC et certains paratextes éditoriaux provenant notamment de Project Gutenberg ou de Wikisource sont retirés en mémoire, sans modification des fichiers sources. Les identifiants des auteurs sont harmonisés et plusieurs contrôles vérifient la correspondance entre les métadonnées et les fichiers, la conservation de chaque œuvre au cours du traitement et l'unicité de son étiquette.

Les romans sont ensuite découpés en segments non chevauchants, en respectant autant que possible les limites de phrases. Trois longueurs cibles sont comparées : 1~000, 1~500 et 2~500 mots. Les phrases sont regroupées jusqu'à atteindre la taille visée ; un dernier fragment inférieur à la moitié de cette taille est fusionné avec le segment précédent. Une phrase exceptionnellement longue peut elle-même être divisée afin d'éviter un écart trop important.

Pour empêcher les œuvres les plus longues de dominer l'apprentissage, chaque roman y fournit au maximum l'équivalent de 50~000 mots. Lorsque ce plafond est dépassé, les segments sont échantillonnés de manière reproductible avec une graine fixée à 42. Tous les segments sont en revanche conservés pour la validation et le test. Avec la longueur de 1~000 mots finalement retenue, l'apprentissage comporte 3~049 segments, la validation 1~080 et le test 1~889. Ces effectifs sont exactement les mêmes dans les conditions ANP et SNP, car les noms propres sont supprimés après la constitution des segments.

\subsection{Deux conditions expérimentales : ANP et SNP}

Dans la condition ANP, aucune suppression particulière des noms propres n'est appliquée. Les noms de personnages, de lieux ou d'organisations peuvent donc entrer dans le vocabulaire du modèle, à condition de satisfaire les filtres statistiques du TF--IDF. Cette version mesure le signal lexical disponible dans son ensemble. Elle peut offrir au classifieur des informations utiles, mais aussi l'inciter à reconnaître des univers fictionnels, des lieux récurrents ou des personnages plutôt que des traits plus généraux.

Dans la condition SNP, chaque segment est analysé avec le modèle français de grande taille de spaCy, \texttt{fr\_core\_}\allowbreak\texttt{news\_lg}, déjà utilisé dans le chapitre précédent. Un token est supprimé lorsqu'il est reconnu comme nom propre par l'étiquette grammaticale \texttt{PROPN}, ou lorsqu'il appartient à une entité de type personne, lieu, organisation ou entité diverse, notées \texttt{PER}, \texttt{LOC}, \texttt{ORG} et \texttt{MISC}. Le nom n'est pas remplacé par une catégorie générique : il est retiré du texte avant la création de la matrice TF--IDF.

Cette suppression est automatique et ne doit pas être considérée comme parfaite. Le modèle linguistique peut manquer certains noms, notamment dans des textes du XIX\textsuperscript{e} siècle, ou supprimer à tort des formes ambiguës. La condition SNP réduit donc la présence d'identifiants directs, mais elle ne neutralise ni le vocabulaire thématique, ni les métiers, ni les objets ou les milieux sociaux représentés. Elle ne produit pas davantage une mesure d'un « style pur » : elle teste seulement la persistance du signal après le retrait d'une catégorie précise d'indices.

\subsection{Représentation des textes et modèles comparés}

Dans les deux conditions, les segments sont représentés par des unigrammes pondérés avec le TF--IDF. Les textes ne sont ni lemmatisés ni filtrés à l'aide d'une liste de mots-outils. Le vectoriseur convertit les formes en minuscules, limite le vocabulaire à 50~000 traits, écarte les termes présents dans moins de deux segments ou dans plus de 80~\% du corpus d'apprentissage, applique une transformation sous-linéaire aux fréquences et normalise chaque vecteur selon la norme L2. Le vocabulaire et les pondérations sont ajustés uniquement sur l'ensemble d'apprentissage, puis appliqués sans réajustement à la validation et au test.

Trois classifieurs sont comparés : un modèle bayésien naïf multinomial, une régression logistique avec pondération équilibrée des classes, et une analyse discriminante linéaire précédée d'une réduction SVD à 100 dimensions puis d'une standardisation. Dans les notebooks, cette dernière méthode est abrégée « LDA », pour \textit{Linear Discriminant Analysis}. Elle est désignée ici par l'acronyme français ADL afin d'éviter toute confusion avec l'allocation de Dirichlet latente présentée dans l'introduction.

Pour chaque condition, la recherche combine les trois tailles de segments, les trois modèles et deux seuils de décision, $0{,}40$ et $0{,}50$, soit dix-huit configurations. Chaque modèle est entraîné sur le seul ensemble d'apprentissage et produit un score d'appartenance à la classe naturaliste pour chaque segment de validation.

\subsection{Sélection des configurations et évaluation}

Les scores des segments d'un même roman sont moyennés afin d'obtenir une décision unique au niveau de l'œuvre. Le critère principal de sélection est le F1 macro calculé séparément pour chacun des trois auteurs de validation, puis moyenné en leur accordant le même poids. La \textit{balanced accuracy} moyenne par auteur, le F1 du moins bon auteur et les mesures globales par œuvre servent ensuite à départager les configurations. Les résultats calculés par segment sont conservés comme diagnostics, mais ne participent pas à la sélection, car plusieurs passages issus d'un même roman ne constituent pas des observations indépendantes.

Cette procédure retient, dans les deux conditions, des segments de 1~000 mots et l'analyse discriminante linéaire. Le seuil choisi sur la validation est de $0{,}50$ pour la condition ANP et de $0{,}40$ pour la condition SNP. Le modèle final est ensuite réentraîné depuis le début sur le seul ensemble d'apprentissage. Les œuvres de validation ne lui sont pas ajoutées, car cela introduirait d'autres auteurs naturalistes dans la classe positive et modifierait la question initiale, fondée sur un apprentissage positif exclusivement zolien.

L'évaluation principale est réalisée au niveau des œuvres, puis répétée séparément pour chaque auteur. Elle repose sur l'exactitude, le F1 macro, la \textit{balanced accuracy}, le rappel de chacune des deux classes et l'aire sous la courbe ROC. Deux analyses de sensibilité examinent également l'effet de la chronologie. La première limite le test aux œuvres publiées entre 1865 et 1903 ; la seconde apparie, chez un même auteur, des œuvres des deux classes publiées à cinq ans d'écart au maximum. L'année n'est jamais fournie au classifieur comme variable prédictive.

Ce protocole réduit plusieurs risques de fuite d'information et permet une comparaison directe entre l'ANP et la SNP. Sa portée reste néanmoins exploratoire. Zola demeure l'unique source positive de l'apprentissage, seulement trois auteurs sont utilisés pour la validation et trois pour le test, et ce dernier ensemble a été consulté au cours du développement. Les résultats ne constituent donc pas une validation définitive d'un classifieur du naturalisme. Ils permettent plutôt d'examiner si le signal lexical appris chez Zola se transfère à quelques trajectoires littéraires distinctes, et dans quelle mesure ce transfert dépend de la présence des noms propres. Les performances des deux conditions sont présentées séparément dans la section suivante avant d'être confrontées.

\section{Résultats des expériences ANP et SNP}
\label{sec:resultats_anp_snp}

Les résultats sont présentés séparément pour chacune des deux conditions. L'évaluation principale porte sur les 25 œuvres du test, tandis que les mesures calculées sur les segments ne sont utilisées que comme diagnostics. Les scores par auteur reçoivent une attention particulière, car les trois écrivains ne sont pas représentés par le même nombre de romans.

\subsection{Résultats de l'expérience ANP}

Sur l'ensemble de validation, la procédure retient l'analyse discriminante linéaire appliquée à des segments de 1~000 mots, avec un seuil de décision fixé à $0{,}50$. Le F1 macro moyen par auteur atteint $0{,}583$ et la \textit{balanced accuracy} moyenne $0{,}675$. Les trois longueurs de segments testées produisent exactement les mêmes résultats principaux au niveau des œuvres ; la taille de 1~000 mots est donc retenue par la règle de départage prévue dans le protocole, et non parce qu'elle serait nettement supérieure. Les résultats restent variables selon les auteurs de validation : le F1 macro vaut $0{,}750$ pour Camille Lemonnier, $0{,}667$ pour Guy de Maupassant et $0{,}333$ pour Georges Eekhoud. Ce dernier résultat ne repose toutefois que sur deux œuvres.

Au niveau des segments du test, l'exactitude atteint $0{,}682$ et le F1 macro $0{,}670$. Le rappel est de $0{,}832$ pour la classe « autre courant » et de $0{,}515$ pour la classe naturaliste. Ces valeurs décrivent le comportement local du modèle, mais elles ne doivent pas être lues comme si les 1~889 segments constituaient autant d'observations indépendantes.

Après agrégation des scores à l'échelle des romans, 18 œuvres sur 25 sont correctement classées. L'exactitude est de $0{,}720$, le F1 macro de $0{,}713$ et la \textit{balanced accuracy} de $0{,}728$. Le modèle reconnaît onze des douze œuvres relevant d'un autre courant et sept des treize œuvres naturalistes. Pour cette dernière classe, la précision atteint $0{,}875$, tandis que le rappel reste limité à $0{,}538$. Ainsi, lorsqu'une œuvre est déclarée naturaliste, cette décision est le plus souvent correcte, mais une partie importante des romans naturalistes ne franchit pas le seuil retenu.

Le tableau~\ref{tab:resultats_anp_auteurs} détaille les résultats pour chacun des trois auteurs du test. L'AUC est calculée à partir des scores continus et ne dépend donc pas du seuil de $0{,}50$.

\begin{table}[H]
    \centering
    \caption{Résultats de l'expérience ANP par auteur du test}
    \label{tab:resultats_anp_auteurs}
    \small
    \setlength{\tabcolsep}{4pt}
    \begin{tabularx}{\textwidth}{>{\raggedright\arraybackslash}X c c c c c c}
        \toprule
        \textbf{Auteur} & \textbf{Correctes} & \textbf{F1} & \textbf{BA}
        & \shortstack{\textbf{Rappel}\\\textbf{nat.}}
        & \shortstack{\textbf{Rappel}\\\textbf{autre}}
        & \textbf{AUC} \\
        \midrule
        Goncourt & 6/10 & 0,583 & 0,750 & 0,500 & 1,000 & 0,813 \\
        J.-K. Huysmans & 7/9 & 0,750 & 0,750 & 0,667 & 0,833 & 0,889 \\
        Octave Mirbeau & 5/6 & 0,778 & 0,750 & 0,500 & 1,000 & 0,875 \\
        \bottomrule
    \end{tabularx}
\end{table}

En accordant le même poids à chaque auteur, le F1 macro moyen s'établit à $0{,}704$, la \textit{balanced accuracy} moyenne à $0{,}750$ et l'AUC moyenne à $0{,}859$. Cette dernière valeur indique que l'ordre des œuvres selon leur score reste informatif, même lorsque le passage à une décision binaire produit plusieurs faux négatifs.

Les sept erreurs permettent de préciser ce résultat. Chez les Goncourt, \textit{Charles Demailly}, \textit{Chérie}, \textit{Madame Gervaisais} et \textit{Manette Salomon} sont étiquetés naturalistes mais classés dans l'autre catégorie. Chez Huysmans, \textit{Marthe, histoire d'une fille} constitue un faux négatif, tandis que \textit{En rade} est le seul faux positif de l'expérience. Enfin, \textit{Le Calvaire} de Mirbeau n'est pas reconnu comme naturaliste. Certaines décisions sont proches du seuil, notamment celles de \textit{Marthe} ($0{,}477$) et de \textit{Madame Gervaisais} ($0{,}465$). D'autres sont beaucoup plus tranchées : \textit{Charles Demailly} et \textit{Chérie} n'obtiennent respectivement que $0{,}072$ et $0{,}189$.

L'examen postérieur des coefficients montre enfin la diversité du signal mobilisé. Parmi les termes fortement associés au prototype zolien figurent des formes discursives comme « lorsque », « tandis », « eut » ou « regardait », des mots liés aux personnes et aux milieux, tels que « enfant », « femme » et « milieu », mais aussi des noms comme « Rougon », « Saccard » ou « Jacques ». À l'inverse, « Bouvard », « Pécuchet », « Arnoux » et « Dambreuse » comptent parmi les termes associés à l'autre classe. Le classement ANP repose donc sur un ensemble mêlant habitudes lexicales, contenus thématiques et éléments propres aux univers romanesques.

\subsection{Résultats de l'expérience SNP}

Dans l'expérience SNP, la validation retient l'analyse discriminante linéaire et des segments de 1~000 mots. Le seuil sélectionné est fixé à $0{,}40$. Le F1 macro moyen par auteur atteint $0{,}637$ et la \textit{balanced accuracy} moyenne $0{,}730$. Les F1 obtenus séparément sont de $0{,}829$ pour Guy de Maupassant, $0{,}750$ pour Camille Lemonnier et $0{,}333$ pour Georges Eekhoud. La sélection demeure donc fragile, notamment parce que le dernier auteur n'est représenté que par deux œuvres.

Le diagnostic au niveau des segments donne une exactitude de $0{,}639$ et un F1 macro de $0{,}600$. Le rappel atteint $0{,}904$ pour la classe « autre courant » et $0{,}343$ pour la classe naturaliste. Ces mesures décrivent les décisions effectuées sur les passages, sans constituer l'unité principale d'interprétation.

Au niveau des œuvres, 17 romans sur 25 sont correctement classés. L'exactitude atteint $0{,}680$, le F1 macro $0{,}653$ et la \textit{balanced accuracy} $0{,}692$. Les douze œuvres de la classe « autre courant » sont reconnues, de même que cinq des treize œuvres naturalistes. La précision de la classe naturaliste est donc de $1{,}000$, pour un rappel de $0{,}385$. Au seuil retenu, le modèle adopte ainsi un comportement très conservateur : aucune œuvre de l'autre classe n'est déclarée naturaliste, mais huit œuvres naturalistes sont laissées de côté.

Le détail par auteur est présenté dans le tableau~\ref{tab:resultats_snp_auteurs}.

\begin{table}[H]
    \centering
    \caption{Résultats de l'expérience SNP par auteur du test}
    \label{tab:resultats_snp_auteurs}
    \small
    \setlength{\tabcolsep}{4pt}
    \begin{tabularx}{\textwidth}{>{\raggedright\arraybackslash}X c c c c c c}
        \toprule
        \textbf{Auteur} & \textbf{Correctes} & \textbf{F1} & \textbf{BA}
        & \shortstack{\textbf{Rappel}\\\textbf{nat.}}
        & \shortstack{\textbf{Rappel}\\\textbf{autre}}
        & \textbf{AUC} \\
        \midrule
        Goncourt & 5/10 & 0,495 & 0,688 & 0,375 & 1,000 & 0,688 \\
        J.-K. Huysmans & 7/9 & 0,679 & 0,667 & 0,333 & 1,000 & 0,944 \\
        Octave Mirbeau & 5/6 & 0,778 & 0,750 & 0,500 & 1,000 & 0,875 \\
        \bottomrule
    \end{tabularx}
\end{table}

Le F1 macro moyen donnant le même poids aux auteurs vaut $0{,}650$, la \textit{balanced accuracy} moyenne $0{,}701$ et l'AUC moyenne $0{,}836$. Le cas de Huysmans illustre la différence entre le classement continu et la décision au seuil : l'AUC atteint $0{,}944$, alors qu'une seule de ses trois œuvres naturalistes est reconnue comme telle. Dans cet échantillon, les scores ordonnent donc assez nettement ses œuvres, mais le seuil de $0{,}40$ ne transforme pas entièrement cet ordre en classifications correctes.

Les huit erreurs sont toutes des faux négatifs. Elles concernent \textit{Charles Demailly}, \textit{Chérie}, \textit{Madame Gervaisais}, \textit{Manette Salomon} et \textit{Renée Mauperin} chez les Goncourt, \textit{En ménage} et \textit{Marthe, histoire d'une fille} chez Huysmans, ainsi que \textit{Le Calvaire} chez Mirbeau. Le score de \textit{La Fille Élisa}, correctement classée, se situe juste au-dessus du seuil à $0{,}414$. La plupart des œuvres mal classées reçoivent en revanche des scores plus éloignés de la frontière de décision ; \textit{Charles Demailly} et \textit{Chérie} obtiennent par exemple $0{,}058$ et $0{,}091$.

Après la suppression automatique des noms propres, les termes les plus associés au prototype zolien comprennent notamment « lorsque », « ça », « nouveau », « frère », « regardait », « milieu », « enfant », « rire », « face » et « femme ». Des termes comme « abbé », « curé » ou « prêtre » contribuent aussi aux scores de certaines œuvres de Mirbeau. Ces coefficients font apparaître à la fois des habitudes lexicales, des verbes de perception ou de parole et des champs thématiques liés à la parenté, au corps, à la religion ou aux milieux sociaux. Leur interprétation reste postérieure à l'évaluation : ils servent à décrire le comportement du modèle, et non à modifier ses paramètres après consultation du test.

\section{Comparaison des expériences ANP et SNP}
\label{sec:comparaison_anp_snp}

Les deux expériences ne se distinguent que par le traitement des noms propres. Elles peuvent donc être mises en regard comme les deux conditions d'une même expérience d'ablation. Le tableau~\ref{tab:comparaison_anp_snp} rassemble les principaux résultats sur le test. La condition ANP obtient les meilleures mesures globales, mais son avantage reste modéré : elle classe correctement 18 œuvres sur 25, contre 17 pour la condition SNP.

\begin{table}[H]
    \centering
    \caption{Comparaison des résultats ANP et SNP sur les œuvres du test}
    \label{tab:comparaison_anp_snp}
    \small
    \setlength{\tabcolsep}{7pt}
    \begin{tabularx}{\textwidth}{>{\raggedright\arraybackslash}X c c}
        \toprule
        \textbf{Indicateur} & \textbf{ANP} & \textbf{SNP} \\
        \midrule
        Œuvres correctement classées & 18/25 & 17/25 \\
        Exactitude & 0,720 & 0,680 \\
        F1 macro par œuvre & 0,713 & 0,653 \\
        \textit{Balanced accuracy} par œuvre & 0,728 & 0,692 \\
        F1 macro moyen par auteur & 0,704 & 0,650 \\
        \textit{Balanced accuracy} moyenne par auteur & 0,750 & 0,701 \\
        AUC moyenne par auteur & 0,859 & 0,836 \\
        Rappel de la classe naturaliste & 0,538 & 0,385 \\
        Rappel de la classe « autre courant » & 0,917 & 1,000 \\
        \bottomrule
    \end{tabularx}
\end{table}

\subsection{Effet du retrait des noms propres}

L'écart le plus visible concerne le rappel de la classe naturaliste. En conservant les noms propres, sept des treize œuvres naturalistes sont reconnues ; après leur suppression, cinq seulement le sont. En contrepartie, le SNP ne produit aucun faux positif : toutes les œuvres qu'il déclare naturalistes le sont selon les étiquettes du corpus. Cette précision parfaite ne doit cependant pas masquer son faible rappel. Elle traduit surtout une décision plus conservatrice, qui laisse huit œuvres naturalistes sous le seuil.

L'avantage global de l'ANP tient à très peu de cas. Les deux conditions donnent la même étiquette à 22 œuvres sur 25. \textit{Renée Mauperin} et \textit{En ménage} sont correctement classées uniquement en ANP, tandis que \textit{En rade} est correctement classé uniquement en SNP. Deux changements favorisent donc l'ANP et un seul le SNP. Sur un échantillon aussi réduit, cette différence ne permet pas d'établir la supériorité générale de la condition avec noms propres.

Cette prudence est renforcée par la validation, où l'ordre était inversé : le F1 moyen par auteur atteignait $0{,}637$ en SNP, contre $0{,}583$ en ANP. Le passage de trois auteurs de validation à trois autres auteurs de test suffit donc à modifier la condition la mieux classée. Ce résultat signale davantage une sélection instable sur de petits effectifs qu'un avantage régulier des noms propres.

Les scores continus conduisent à la même prudence. Leur corrélation de rang atteint $0{,}972$ : les deux modèles ordonnent presque toujours les romans de la même manière. Tous les scores SNP sont néanmoins inférieurs à leur équivalent ANP, leur moyenne passant de $0{,}379$ à $0{,}235$. Ce déplacement ne mesure pas directement la « part » des noms propres, car les modèles sont réentraînés dans deux espaces lexicaux différents et leurs seuils ont été réglés séparément, à $0{,}50$ et $0{,}40$. Il montre plutôt que la suppression modifie l'échelle des scores et réduit la propension du modèle à attribuer la classe naturaliste, sans bouleverser l'ordre général des œuvres.

L'effet varie d'ailleurs selon les auteurs. Les décisions concernant Mirbeau sont identiques dans les deux expériences. Chez Huysmans, l'AUC progresse en SNP alors que le F1 diminue, ce qui indique un classement continu pertinent mais mal traduit par le seuil. La baisse moyenne provient surtout des Goncourt. Le retrait des noms propres n'a donc pas un effet uniforme sur les différents univers littéraires.

\subsection{Nature du signal et effet de la chronologie}

Le maintien d'une AUC moyenne de $0{,}836$ en SNP constitue un résultat important : les personnages et les lieux ne suffisent pas à expliquer le classement. Un signal lexical appris à partir de Zola subsiste chez des auteurs absents de l'apprentissage. Il ne correspond pas pour autant à un style naturaliste isolé. Les termes encore disponibles associent des formes discursives, comme « lorsque » ou « regardait », à des mots liés à la parenté, au corps, à la religion et aux milieux sociaux. Avec des unigrammes TF--IDF, habitudes d'écriture et contenu des romans restent ainsi mêlés.

La chronologie représente un autre facteur de confusion. Lorsque le test est limité aux 21 œuvres publiées entre 1865 et 1903, les deux conditions atteignent la même exactitude de $0{,}714$. Les F1 macro restent proches, avec $0{,}697$ en ANP et $0{,}679$ en SNP. La petite avance globale de l'ANP n'est donc pas robuste à cette restriction. Celle-ci conserve toutefois des écarts de dates importants à l'intérieur de la période et ne neutralise pas l'évolution propre à chaque écrivain. L'appariement d'œuvres des deux classes publiées à cinq ans d'écart au maximum va dans le même sens : l'ANP ordonne correctement trois des quatre paires et le SNP deux. Quatre paires sont trop peu nombreuses pour lever le doute sur un effet de période.

\subsection{Portée, limites et prolongements}

Le protocole présente plusieurs points forts. Les auteurs de l'apprentissage, de la validation et du test sont disjoints, les copies exactes sont recherchées, et l'évaluation principale porte sur les œuvres plutôt que sur des segments corrélés. Le calcul séparé par auteur évite aussi que les Goncourt, plus représentés, déterminent seuls le résultat. Enfin, l'ablation ANP/SNP porte sur les mêmes textes et montre que le signal ne disparaît pas avec les identifiants fictionnels. Les deux modèles font nettement mieux qu'une décision fondée uniquement sur la classe majoritaire du test, qui ne classerait correctement que 13 œuvres sur 25.

Ces garanties ne suffisent cependant pas à faire de l'expérience une validation définitive. Zola demeure l'unique auteur positif de l'apprentissage : le modèle apprend nécessairement un mélange de naturalisme, le lexique zolien et de contenus propres à ses romans. Le test ne réunit que trois auteurs et 25 œuvres, et il a été consulté pendant le développement. Les catégories littéraires peuvent en outre être discutées pour certaines œuvres de transition. Enfin, la suppression automatique des noms propres n'est pas infaillible et son taux d'erreur n'a pas été mesuré sur ce corpus du XIX\textsuperscript{e} siècle.

Plusieurs tests complémentaires permettraient d'approfondir ce travail. La priorité serait de constituer un nouveau test réellement aveugle, avec davantage d'auteurs et des œuvres mieux équilibrées par période. L'apprentissage pourrait ensuite intégrer plusieurs écrivains naturalistes, puis évaluer successivement un auteur tenu à l'écart ; une telle validation distinguerait mieux le courant littéraire de la seule signature de Zola. Un modèle fondé uniquement sur l'année fournirait par ailleurs un témoin simple, à dépasser avant d'attribuer les résultats au texte. Enfin, la comparaison de familles de traits --- mots-outils, ponctuation, n-grammes de caractères, catégories grammaticales ou relations syntaxiques --- aiderait à séparer les régularités formelles du vocabulaire thématique. Pour la condition SNP, remplacer les entités par des catégories génériques, comme \texttt{[PER]} ou \texttt{[LOC]}, et vérifier manuellement un échantillon permettrait aussi de conserver la structure des phrases tout en contrôlant la qualité de la reconnaissance.

\section*{Conclusion du chapitre}
\addcontentsline{toc}{section}{Conclusion du chapitre}

Ce chapitre montre qu'une partie du signal appris à partir de Zola se transfère à des œuvres naturalistes de Goncourt, Huysmans et Mirbeau, alors même que ces auteurs sont absents de l'apprentissage. Les résultats sont supérieurs à une règle majoritaire et le classement des œuvres reste largement informatif. Surtout, la forte proximité entre les scores ANP et SNP indique que ce transfert ne repose pas uniquement sur les noms de personnages ou de lieux.

La réponse à la problématique demeure toutefois partielle. Le modèle met en évidence une proximité textuelle avec Zola ; il ne démontre pas encore l'existence de marqueurs du naturalisme indépendants de l'auteur, des thèmes et de la chronologie. Le faible rappel des œuvres naturalistes, la variation entre auteurs et la fragilité des contrôles temporels invitent à ne pas transformer ces résultats en critères de classement littéraire. Ils restent néanmoins exploitables dans le cadre d'un test exploratoire : ils valident l'intérêt du protocole, rendent visibles ses principaux biais et fournissent une base précise pour une expérimentation ultérieure, plus large et véritablement confirmatoire.


\nocite{*}
\printbibliography[nottype=online, heading=bibintoc, title={Bibliographie}]
\printbibliography[type=online, heading=bibintoc, title={Sitographie}]


\appendix
\chapter*{Annexe technique }
\addcontentsline{toc}{chapter}{Annexe }

\section{Reflexion sur le stockage des données}

Après avoir recupéré un nombre consequant de romans naturalistes et non naturaliste, il a fallu réflechir à un moyen de les stocker quelque part. Pour cela, j'ai réfléchi à plusieurs idées pour procécéder. Le choix c'est porté sur une base de données SQL avec comme interface PHP admin. Cette base de données à d'abord été remplie en local pour effectuer des tests et ensuite, une fois que tout fonctionnait correctement, j'ai pu extraire un fichier \textsl{Dumps} et l'importer dans PHP Admin.

\section{Nettoyage des textes}

Après avoir réfléchi à la manière de stocker les romans, il a fallu réfléchir à un moyen de les nettoyer. En effet, les romans contiennent beaucoup de caractères spéciaux qui ne sont pas utiles pour l'analyse, notamment les caractères de mise en forme, la pagination, les notes de bas de page, les chapitres ainsi que les prologues. En effet, les prologues sont souvent des textes qui ne sont pas écrits par l'auteur du roman et qui ne sont donc pas représentatifs de son style d'écriture.

Les romans naturalistes et non naturalistes sont issus de types de données différents. Il y a trois groupes de formats de données : les romans au format \textsl{.epub}, les romans au format \textsl{.pdf} et les romans au format \textsl{.txt}\footnote{Ici, les textes au format \textsl{.txt} sont organisés en dossiers, c'est-à-dire qu'un fichier \textsl{.txt} correspond à une page du roman. Il a donc fallu les regrouper pour obtenir un seul fichier \textsl{.txt} par roman.}.

Pour ce qui est des textes au format \textsl{.epub}, j'ai utilisé le logiciel \textsl{Calibre}\footnote{Lien du logiciel : \href{https://calibre-ebook.com/fr/about}{https://calibre-ebook.com/fr/about}, site consulté le 18 juin 2026.} pour les convertir au format \textsl{.txt}. Pour les textes au format \textsl{.pdf}, j'ai utilisé le logiciel \textsl{Adobe Acrobat Pro} afin de les convertir au format \textsl{.txt}. Pour les textes au format \textsl{.txt}, j'ai utilisé un script Python\footnote{Le script est disponible en annexe.} pour regrouper les fichiers en un seul fichier par roman.

La grande majorité de mes romans étaient au format \textsl{.epub}. Après leur conversion au format \textsl{.txt}, j'ai donc utilisé un script Python\footnote{Le script est disponible en annexe. Il faut préciser que je n'ai pas utilisé un seul script, mais plusieurs. Il s'agira donc d'un script général. Il y a eu énormément de modifications internes, notamment pour les chapitres qui sont soit en majuscules, avec chiffres romains, ou sans chiffres romains. Il a donc fallu faire des ajustements à chaque fois.} pour nettoyer les textes. Ce script supprime les caractères spéciaux, les notes de bas de page, les chapitres et les prologues. Il supprime également les lignes vides et les espaces en trop. Le script est basé sur des expressions régulières pour identifier les éléments à supprimer.

Pour ce qui est du choix des éléments à supprimer, j'ai fait le choix de supprimer les chapitres, car ils ne sont pas représentatifs du style d'écriture de l'auteur. De plus, comme nous l'avons vu dans le chapitre 2, j'ai passé un détecteur de motifs sur les textes afin d'identifier les motifs récurrents. J'ai donc gardé uniquement le corps du texte pour l'analyse. Ce choix se justifie également par la volonté d'obtenir une normalisation entre les textes\footnote{certains textes de disposant pas de }(une ligne est égal a une phrase), ce choix peut etre mit en opposition avec le choix de garder la distinction entre les différents chapitres du romans mais je me defend en prenant en comptre que certains romans ne sont pas découpés en chapitres. Les differents romans sont soient découpés en partie (exemple : Partie I, Partie II, etc.), soit en chapitres (exemple : Chapitre I, Chapitre II, etc.), soit avec une combinaison des deux (exemple : Partie I, Chapitre I, Chapitre II, Partie II, Chapitre III, etc.). Ou pour finir, le titre du chapitre est écrit en majuscules, avec des chiffres romains, ou sans chiffres romains. Il aurait donc fallu faire des ajustements à chaque fois pour identifier les chapitres. De plus, certains romans résultaient de l'assemblage de plusieurs fichiers \textsl{.txt}. Dans ces cas, la séparation entre les chapitres, initialement présente dans la structure HTML du roman, n'était plus conservée.

\section{Création de la base de données}

Après avoir nettoyé les textes, il a fallu créer une base de données pour stocker les informations relatives aux romans. J'ai choisi d'utiliser le système de gestion de base de données relationnelle \textsl{MySQL}, car il permet d'organiser les données sous forme de tables et de conserver des relations claires entre les auteurs, les romans et les unités textuelles extraites.

Je vais donc présenter la structure de cette base de données ainsi que les différentes tables qui la composent. Chaque table a un rôle précis : certaines servent à stocker les informations bibliographiques, tandis que d'autres permettent de conserver les textes, les segments ou encore les phrases extraites pour l'analyse.

J'ai choisi de présenter ce code directement dans le corps du mémoire\footnote{Je précise également que tous les codes utilisés pour la réalisation de ce mémoire sont disponibles sur mon GitHub.} et non uniquement en annexe, afin de rendre plus claire la manière dont les données ont été organisées avant l'analyse. 

\begin{listing}[htbp]
\begin{minted}[
    fontsize=\scriptsize,
    breaklines=true,
    frame=lines,
    linenos=true,
    numbersep=5pt,
    autogobble
]{sql}
CREATE DATABASE romans_nlp;
USE romans_nlp;

-- Creation des tables pour la base de donnees des romans et des auteurs
CREATE TABLE auteurs (
    id_auteur INT PRIMARY KEY AUTO_INCREMENT,
    nom VARCHAR(100) NOT NULL,
    prenom VARCHAR(100)
);

-- Table pour stocker les romans, avec leur titre et leur annee de publication
CREATE TABLE romans (
    id_roman INT PRIMARY KEY AUTO_INCREMENT,
    titre VARCHAR(255) NOT NULL,
    annee_publication INT,
    texte_brut LONGTEXT
);

-- Table d'association pour gerer la relation plusieurs-a-plusieurs entre auteurs et romans
CREATE TABLE auteurs_romans (
    id_auteur INT NOT NULL,
    id_roman INT NOT NULL,
    PRIMARY KEY (id_auteur, id_roman),
    FOREIGN KEY (id_auteur) REFERENCES auteurs(id_auteur),
    FOREIGN KEY (id_roman) REFERENCES romans(id_roman)
);

-- Table pour stocker les segments extraits des romans, avec leur ordre de segmentation
CREATE TABLE segments (
    id_segment INT PRIMARY KEY AUTO_INCREMENT,
    id_roman INT NOT NULL,
    ordre_segment INT NOT NULL,
    texte_segment LONGTEXT NOT NULL,
    FOREIGN KEY (id_roman) REFERENCES romans(id_roman)
);

/*
Apres reflexion, il est egalement possible d'ajouter les dates
de naissance et de mort des auteurs afin d'enrichir la base de donnees.
*/
ALTER TABLE auteurs
ADD date_naissance INT;

ALTER TABLE auteurs
ADD date_mort INT;

-- Table pour stocker les phrases extraites des romans, avec leur motif d'extraction
CREATE TABLE phrase_roman (
    id_phrase INT PRIMARY KEY AUTO_INCREMENT,
    id_roman INT NOT NULL,
    ordre_phrase INT NOT NULL,
    texte_phrase TEXT NOT NULL,
    motif_phrase TEXT NOT NULL,
    FOREIGN KEY (id_roman) REFERENCES romans(id_roman)
);
\end{minted}
\caption{Création de la base de données romans\_nlp}
\end{listing}

Ce code permet de créer la base de données utilisée pour stocker les romans du corpus ainsi que les informations nécessaires à leur analyse. La base, nommée \texttt{romans\_nlp}, est d'abord créée puis sélectionnée afin que toutes les tables suivantes y soient enregistrées.

La première table, \texttt{auteurs}, sert à stocker les informations relatives aux auteurs. Chaque auteur reçoit un identifiant unique, auquel sont associés son nom et, lorsque l'information est disponible, son prénom. Cette table est ensuite complétée par deux colonnes supplémentaires, \texttt{date\_naissance} et \texttt{date\_mort}, afin d'enrichir la description des auteurs et de permettre d'éventuelles analyses chronologiques.

La table \texttt{romans} contient les informations principales sur les œuvres du corpus. Elle associe à chaque roman un identifiant unique, un titre, une année de publication ainsi que le texte brut du roman. Le choix du type \texttt{LONGTEXT} pour le champ \texttt{texte\_brut} permet de stocker des textes longs, ce qui est nécessaire dans le cas de romans complets.

La table \texttt{auteurs\_romans} permet de faire le lien entre les auteurs et les romans. Elle sert à gérer une relation de type plusieurs-à-plusieurs : un auteur peut être associé à plusieurs romans, et un roman peut éventuellement être associé à plusieurs auteurs. Les clés étrangères \texttt{id\_auteur} et \texttt{id\_roman} permettent de relier cette table aux tables \texttt{auteurs} et \texttt{romans}.

La table \texttt{segments} sert à stocker les segments extraits des romans. Chaque segment est rattaché à un roman grâce à la clé étrangère \texttt{id\_roman}. Le champ \texttt{ordre\_segment} permet de conserver l'ordre d'apparition des segments dans le texte d'origine. Cela est important, car même si le texte est découpé pour les traitements, il reste possible de retrouver sa progression initiale.

Enfin, la table \texttt{phrase\_roman} permet de stocker les phrases extraites des romans. Chaque phrase est associée au roman dont elle provient, à son ordre d'apparition, à son contenu textuel et au motif qui a permis son extraction. Ce dernier champ est important, car il permet de garder une trace de la méthode utilisée pour sélectionner les phrases et de relier chaque extraction à un critère précis.

L'ensemble de cette structure permet donc de conserver à la fois les textes complets, leurs découpages en segments, les phrases extraites et les informations bibliographiques liées aux auteurs et aux œuvres. Elle constitue ainsi une base organisée pour les traitements de fouille de texte et d'analyse automatique du corpus.


\section{Création des informations pour la base de données}

Après avoir créé les tables nécessaires pour stocker les informations sur les auteurs, les romans et les relations entre eux, il est temps d'insérer les données relatives aux frères Goncourt et à leurs œuvres. Le code suivant présente donc l'insertion des auteurs, des romans et des relations entre ces deux tables.

\begin{minted}[
    fontsize=\scriptsize,
    breaklines=true,
    frame=lines,
    linenos=true,
    numbersep=5pt,
    autogobble
]{sql}
/*
Insertion des freres Goncourt.
*/
INSERT INTO auteurs (nom, prenom) 
VALUES
    ('de Goncourt', 'Edmond'),
    ('de Goncourt', 'Jules');
/*
Insertion des romans des freres Goncourt
*/
INSERT INTO romans (titre, annee_publication) 
VALUES
    ('Charles Demailly', 1860),
    ('Soeur Philomene', 1861),
    ('Renee Mauperin', 1864),
    ('Germinie Lacerteux', 1865),
    ('Manette Salomon', 1867),
    ('Madame Gervaisais', 1869),
    ('La Fille Elisa', 1877),
    ('Freres Zemganno', 1879),
    ('La Faustin', 1882),
    ('Cherie', 1884);
/*
Insertion des relations entre les auteurs et les romans.
Les six premiers romans sont associes aux deux freres Goncourt.
Les romans suivants sont associes uniquement a Edmond de Goncourt.
*/
INSERT INTO auteurs_romans (id_auteur, id_roman) 
VALUES
    (1,1), (2,1),
    (1,2), (2,2),
    (1,3), (2,3),
    (1,4), (2,4),
    (1,5), (2,5),
    (1,6), (2,6),
    (1,7),
    (1,8),
    (1,9),
    (1,10);
-- Mise a jour des dates de naissance et de mort de Jules de Goncourt.
UPDATE auteurs
SET date_naissance = 1830,
    date_mort = 1870
WHERE id_auteur = 2;

-- Mise a jour des dates de naissance et de mort d'Edmond de Goncourt.
UPDATE auteurs
SET date_naissance = 1822,
    date_mort = 1896
WHERE id_auteur = 1;

\end{minted}

Ce code permet d'insérer dans la base de données les informations relatives aux frères Goncourt et à leurs romans. Dans un premier temps, les deux auteurs sont ajoutés dans la table \texttt{auteurs}. Ils reçoivent chacun un identifiant, qui permettra ensuite de les relier aux œuvres présentes dans la base.

La deuxième partie du code insère les romans des frères Goncourt dans la table \texttt{romans}. Chaque œuvre est associée à son titre et à son année de publication. Cette étape permet donc de constituer une partie du corpus à partir des romans retenus pour l'analyse.

Ensuite, la table \texttt{auteurs\_romans} est utilisée pour associer les auteurs aux romans. Cette table est importante, car elle permet de gérer le cas particulier des œuvres écrites à deux mains. Les six premiers romans sont associés à Edmond et Jules de Goncourt, tandis que les romans publiés après la mort de Jules sont associés uniquement à Edmond. Cela permet de conserver une information plus précise sur l'attribution des œuvres.

Enfin, les deux dernières requêtes mettent à jour les dates de naissance et de mort des auteurs. Ces informations biographiques ne sont pas directement nécessaires pour stocker les textes, mais elles permettent d'enrichir la base de données. Elles peuvent aussi être utiles par la suite si l'on souhaite croiser les résultats de l'analyse avec des éléments chronologiques ou biographiques.

Ce passage sert donc à remplir une première partie de la base de données avec les auteurs, les romans et les relations entre eux. Il montre aussi l'intérêt de séparer les informations dans plusieurs tables : les auteurs sont stockés d'un côté, les romans de l'autre, et la relation entre les deux est conservée dans une table spécifique.

Il faut maintenant ajouter les romans dans la base de données, il peut y avoir plusieurs façons de le faire, 

\end{document}